\chapter{Reference: \code{live} (inline loop)}
\label{ch:refman-live}

\begin{aside}
\code{maya::live} is \code{run} without the alt-screen.
You pass a lambda that returns an Element; the runtime calls
it at the target frame rate; the rendered frame replaces
the previous one in place, in scrollback.
\end{aside}

\section{Signature}

\begin{lstlisting}[language=C++]
namespace maya {
    struct LiveConfig { int fps = 30; };
    void live(LiveConfig cfg, std::function<Element()> view);
    void quit();
}
\end{lstlisting}

\section{Working example}

\begin{lstlisting}[language=C++]
#include <maya/maya.hpp>
#include <chrono>

using namespace maya;
using namespace maya::dsl;

int main() {
    auto start = std::chrono::steady_clock::now();

    live({.fps = 30}, [&]() -> Element {
        auto elapsed =
            std::chrono::duration_cast<std::chrono::milliseconds>(
                std::chrono::steady_clock::now() - start);

        if (elapsed >= std::chrono::seconds(5)) {
            quit();
        }

        float frac = std::min(1.0,
            elapsed.count() / 5000.0);
        int filled = static_cast<int>(frac * 30);
        std::string bar;
        for (int i = 0; i < 30; ++i) {
            bar += (i < filled) ? "#" : ".";
        }
        return h(
            text(bar) | Fg<100, 220, 160>,
            text(" "),
            text(std::format("{:>3}\%",
                static_cast<int>(frac * 100))) | Bold
        );
    });

    print(t<"done"> | Bold | Fg<100, 255, 140>);
    return 0;
}
\end{lstlisting}

\section{Notes}

\begin{itemize}[leftmargin=*]
  \item Call \code{quit()} from inside the lambda to break
    the loop after this frame.
  \item Output stays in scrollback; redirecting to a file
    captures the last frame.
  \item The lambda runs on the main thread; do not block
    inside it.
\end{itemize}

\section{Recap}

\begin{recap}
\begin{itemize}[leftmargin=*]
  \item Lambda returns an Element each frame.
  \item Inline; no alt-screen.
  \item \code{quit()} ends the loop.
\end{itemize}
\end{recap}

\section*{Workshop}

\begin{enumerate}[leftmargin=*]
  \item Replace a \code{while(true) sleep} progress loop
    with \code{live}.
  \item Show a spinner during a long subprocess call.
  \item Render a small dashboard for the duration of a
    deploy.
\end{enumerate}
